\section{Banks and the second round}
\label{sec:banks}

This module tests the last clause of the central claim, that the first-round burden does not fall
on aggregate bank capital. Displacement reaches banks, but mostly through what lost spending does
to everyone else rather than through the borrowers who lose their jobs. We apply the loss across
\result{NInstitutions} individual balance sheets and ask which are short of capital afterwards,
and then whether household relief is the instrument that protects them. The balance sheets are
measured; every loss applied above the level inside the data is scenario.

\subsection{Almost nothing happens at the level inside the data}

At ten percent of the wage bill displaced, system losses are about \result{SysLossTen} billion
dollars, of which \result{SysFirstTen} billion is first round and \result{SysSecondTen} billion,
\result{SecondRoundShareTen} percent of the total, arrives through the second round.
\result{BreachTenN} institutions fall below their capital threshold, holding
\result{BreachTenAssets} percent of system assets, against a baseline of
\result{BaselineBreaches} institutions and \result{BaselinePctAssets} percent before any loss.
That is a banking system that barely moves.

Breaches open above the level inside the data: \result{BreachTwentyFiveN} institutions holding
\result{BreachTwentyFiveAssets} percent of assets at twenty-five percent displacement, and
\result{BreachFiftyN} holding \result{BreachFiftyAssets} percent at fifty. Allowing one year of
earnings to absorb losses before the capital test cuts those to
\result{BreachTwentyFiveAssetsEarnings} and \result{BreachFiftyAssetsEarnings} percent, the single
largest sensitivity in this module and the reason the upper rows are shape rather than level.

\subsection{Which institutions}

The result worth reporting is the dispersion, in Figure~\ref{fig:breach}. Card-heavy lenders are
the most exposed by a distance, \result{CardHeavyTwentyFiveInst} percent of them breaching at
twenty-five percent displacement and holding \result{CardHeavyTwentyFiveAssets} percent of that
class's assets, rising to \result{CardHeavyFiftyInst} percent of institutions at fifty. Credit
unions are the only class with any stress at the ten percent level,
\result{CreditUnionTenInst} percent of them, because their book is almost entirely household
credit with no commercial business to dilute it. Mortgage portfolio lenders are among the least
exposed at every level, \result{MortgageLenderTwentyFiveInst} percent of institutions at
twenty-five percent.

That last result settles at institution level a question occupational data alone leave open. What
makes a lender fragile to a wage shock is an unsecured consumer book, whose loss given default is
close to one, rather than a mortgage book concentrated in one region or industry. Two
qualifications travel with it: losses are allocated pro rata by category, so business mix is the
only source of dispersion here and the true spread is wider; and the largest institutions are
diversified by construction, so this is a statement about the tail rather than the center.

\begin{figure}[htbp]
\centering
\includegraphics[width=0.84\textwidth]{fig5_breach.pdf}
\caption{Share of each business model's assets held by institutions below their capital threshold,
at three displacement levels. Balance sheets measured at 30 June 2026; loss application is
scenario, and the twenty-five and fifty percent levels are outside the observed data.}
\label{fig:breach}
\end{figure}

\subsection{The instrument test: household relief does not protect bank capital}

If the second round is what reaches banks, instruments acting on the first round should not do
much for bank capital. Table~\ref{tab:relief} tests that through the same balance sheets.
Forbearance, automatic income-driven student repayment and enhanced wage insurance together
remove \result{ReliefRemovedTen} percent of system bank losses at the ten percent level,
\result{ReliefRemovedTwentyFive} percent at twenty-five and \result{ReliefRemovedFifty} percent at
fifty: under a tenth everywhere. The arithmetic is in the first paragraph of this section, since
only \result{SysFirstTen} billion of a \result{SysLossTen} billion total is first round. The most
expensive instrument in the set costs \result{WageInsuranceCost} billion dollars and removes
\result{WageInsuranceRemoved} billion of bank losses.

\begin{table}[htbp]
\centering
\caption{Relief instruments against system bank losses at ten percent of the wage bill displaced.
Cells with no sourced fiscal cost read n/a. Scenario.}
\label{tab:relief}
\begin{tabular}{lrrr}
\toprule
Instrument & System loss removed (USD bn) & Per cent of system loss & Fiscal cost (USD bn) \\
\midrule
mortgage forbearance, partial & 2.60 & 1.27 & n/a \\
auto and consumer forbearance, partial & 5.35 & 2.61 & n/a \\
mortgage forbearance, full & 4.42 & 2.16 & n/a \\
auto and consumer forbearance, full & 9.10 & 4.44 & n/a \\
income-driven student repayment, automatic enrolment & 7.55 & 3.68 & 7.3 \\
wage insurance, current US, 0.50 for 26 weeks & 4.60 & 2.24 & 112.2 \\
wage insurance, enhanced, 0.70 for 52 weeks & 12.87 & 6.28 & 314.3 \\
All together: forbearance & 19.71 & 9.61 & 317.6 \\
\bottomrule
\end{tabular}

\end{table}

The conclusion is about what these instruments are for.

\begin{quote}
Household relief must be judged on household outcomes, not on bank capital. It does reshape the
distribution of failures: at the twenty-five percent level it takes assets in breach from
\result{ReliefAssetsBeforeTwentyFive} to \result{ReliefAssetsAfterTwentyFive} percent and stops
\result{ReliefStoppedTwentyFive} institutions breaching. It changes who fails, not how much is
lost.
\end{quote}

Those two facts are consistent. The instruments remove a small share of a large aggregate loss,
but they remove it from the institutions whose books are concentrated in the relieved categories,
so a small total with concentrated incidence changes many capital outcomes while barely moving the
system total. An instrument chosen to protect capital would have to act on the second round, which
means acting on spending, house prices and business credit.

\subsection{The exposures this engine does not reach}

Four sets of institutions are exposed in ways the panel does not capture, and naming them is more
useful than putting a number on them the data cannot support. \emph{State and local government} is
exposed through its own tax base, \result{StateLocalWageLinkedBn} billion dollars of wage-linked
income tax, \result{StateLocalWageLinkedShare} percent of its own tax receipts, and cannot run the
deficit through because of balanced-budget requirements; its in-year spending cut feeds the same
demand channel that produces nine tenths of the bank losses above, and that loop is not in the
engine, so its absence makes our figures conservative in a direction we can sign but not size.
\emph{Pension funds} are exposed on both sides at once, benefits fixed in nominal terms while
contributions are a share of covered payroll, so displacement widens the funding gap mechanically
and raises the required contribution rate on the employers who remain; the accounts capture only
the \result{PensionHoldingsBn} billion dollars they hold. \emph{Insurers and private credit} hold
both sides, insurers \result{InsurerHoldingsBn} billion dollars of wage-backed claims against
long-dated liabilities, while private credit cannot be separated from the other-financial
aggregate, so the \result{OtherFinancialHoldingsBn} billion attributed to that sector is an upper
bound by an unknown margin. And everything outside the insured depository system, nonbank
servicers and consumer lenders among them, is absent from the panel; as a class they are more
thinly capitalized than the institutions we measure, so the system figures are again a lower
bound.
